\documentclass[10pt]{article}
\usepackage{amsmath,amssymb,geometry}
\geometry{margin=1in}
\setlength{\emergencystretch}{3em}
\title{A counterexample to conjecture 00000008438}
\author{gaochengzhecpu (AI-assisted submission)}
\date{3 October 2026}
\usepackage{url}
% Compact consistent title for a short mathematical note.
\makeatletter
\renewcommand{\maketitle}{\begin{center}
{\Large\bfseries\@title\par}\vspace{0.6em}
{\small\@author\par\@date}
\end{center}\vspace{0.5em}}
\makeatother
\setlength{\emergencystretch}{3em}
\begin{document}
\maketitle
\paragraph{Claim being disproved.}
The conjecture bounds the number of columns of a strength-$t$ orthogonal array over $n$ symbols by $n^2-n+t$. Here ``strength $t$'' means that on any $t$ distinct rows every ordered $t$-tuple of symbols occurs equally often. We use an array of index one, so the example does not rely on repeating columns.
\paragraph{Construction.}
Take $n=2$, $t=3$, and the array
\[
 A=\begin{pmatrix}
 0&1&0&1&0&1&0&1\\
 0&0&1&1&0&0&1&1\\
 0&0&0&0&1&1&1&1
 \end{pmatrix}.
\]
Its columns are all eight distinct binary words of length three. Hence any ordering of its three rows contains every ordered binary triple exactly once: reordering coordinates is a bijection on $\{0,1\}^3$. Thus $A$ is a strength-three, index-one orthogonal array.
\paragraph{Contradiction.}
The array has $8$ columns, but the proposed bound is
\[
 n^2-n+t=2^2-2+3=5<8.
\]
A counterexample to the first bound disproves the conjunction in the conjecture; the prime-power attainability and deficit assertions need not be considered.
\paragraph{A second witness with matching row and alphabet counts.}
Take instead the three-row array whose columns are all $27$ words in $\{0,1,2\}^3$. It also has strength three and index one, since permuting its rows permutes all ternary triples. Here both the alphabet size and the row count equal $3$, while $27>3^2-3+3=9$. Thus this example also avoids any possible ambiguity between those two meanings of $n$. The accompanying Lean theorems \path{ternaryArray_is_OA3}, \path{ternaryArray_columns_distinct}, and \path{conjecture8438_false_ternary} verify the full second witness and the resulting negation.
\paragraph{Formal verification and semantic correspondence.}
The accompanying \texttt{lean/Main.lean} uses only Lean 4.19.0 and \texttt{Std}. The type \texttt{Fin r} indexes rows, \texttt{Fin c} indexes columns, and \texttt{Fin n} indexes symbols. The definition \texttt{occurrences} counts actual columns matching prescribed entries. The predicate \texttt{OA3} quantifies over every triple of distinct rows, in every order, and every prescribed symbol triple. The theorem \texttt{binaryArray\_is\_OA3} verifies this property for the displayed array; \path{binaryArray_columns_distinct} also verifies that no two columns coincide.

\texttt{ColumnBoundAtThree} states the conjectured universal upper bound specialized to $t=3$, over arbitrary alphabet size, row count, column count, and array. The final theorem \texttt{conjecture8438\_false} proves its negation by instantiating that universal claim with $A$. Thus the formal result includes the orthogonal-array hypotheses and the implication to the claimed bound, rather than only the numerical inequality $8>5$.

All enumeration is checked by Lean's kernel using \texttt{decide}; there is no external computation oracle. The printed axiom lists contain only the standard Lean axiom \texttt{propext}. No \texttt{sorry}, \texttt{admit}, \texttt{native\_decide}, or additional axiom is used.
\paragraph{Reproduction.}
Run \texttt{lean Main.lean} with Lean 4.19.0. Successful compilation prints the axiom lists of the array property, column distinctness, and final disproof.
\end{document}

